% Part: first-order-logic
% Chapter: models-theories
% Section: expressing-props-of-structures

\documentclass[../../../include/open-logic-section]{subfiles}

\begin{document}

\olfileid{fol}{mat}{exs}

\olsection{\printtoken{P}{structure}ల ధర్మాలను వ్యక్తం చేయడం}

\begin{explain}
ప్రమేయాలు, సంబంధాలపై షరతులను వ్యక్తం చేయడం, లేదా మరింత సాధారణంగా ఒక
నిర్మాణంలోని ప్రమేయాలు, సంబంధాలు ఆ షరతులను సంతృప్తిపరుస్తాయని వ్యక్తం
చేయడం తరచుగా ఉపయోగకరమూ ముఖ్యమూ అవుతుంది. ఉదాహరణకు, ఒక భాష కోసం మనం
\tetoken{విధేయ సంకేతాలకు}{!!{predicate}s} ఉద్దేశించిన ``అర్థాన్ని'' ``పట్టుకునే'' \tetoken{నిర్మాణాలను}{!!{structure}s},
అలా చేయని వాటి నుంచి వేరు చేయడానికి మార్గాలు కావాలి. ఏ \tetoken{నిర్మాణాలను}{!!{structure}s}
``ఉద్దేశించాలో'' నిర్దేశించడంలో మనకు పూర్తి స్వేచ్ఛ ఉంది. ఉదాహరణకు,
\tetoken{విధేయ సంకేతం}{!!{predicate}}~$\le$ యొక్క అర్థనిర్దేశం ఒక క్రమం కావాలని, లేదా క్షేత్రం
సమితులతో ఏర్పడి $\Obj \in$కు ``దీనిలో ఒక \tetoken{మూలకం}{!!a{element}}'' అనే సంబంధమే
అర్థనిర్దేశంగా ఉండే~$\Lang L$ అర్థనిర్దేశాలపైనే మనకు ఆసక్తి ఉందని
నిర్దేశించవచ్చు. కానీ భాషలోని \tetoken{వాక్యాలతో}{!!{sentence}s} ఇలా చేయగలమా? మరోలా
చెప్పాలంటే, \tetoken{ఒక నిర్మాణం}{!!a{structure}}~$\Struct M$పై ఏ షరతులను~$\Struct M$ భాషలోని
\tetoken{ఒక వాక్యంతో}{!!a{sentence}} (లేదా బహుశా \tetoken{వాక్యాలు}{!!{sentence}s} సమితితో) వ్యక్తం చేయగలం?
మనం వ్యక్తం చేయలేని కొన్ని షరతులు ఉన్నాయి. ఉదాహరణకు,
$\Domain M = \Nat$ అయిన ఒక \tetoken{నిర్మాణం}{!!{structure}}~$\Struct M$లో మాత్రమే సత్యమయ్యే
$\Lang L_A$ వాక్యం ఏదీ లేదు. ``క్షేత్రంలో సహజ సంఖ్యలు మాత్రమే ఉన్నాయి''
అని మనం వ్యక్తం చేయలేం. కానీ \tetoken{నిర్మాణాలకు}{!!{structure}s} చెందిన కొన్ని
``నిర్మాణాత్మక ధర్మాలను'' బహుశా వ్యక్తం చేయగలం. \tetoken{నిర్మాణాలు}{!!{structure}s} యొక్క ఏ
ధర్మాలను \tetoken{వాక్యాలు}{!!{sentence}s} ద్వారా వ్యక్తం చేయగలం? మరో విధంగా అడిగితే,
\tetoken{ఒక వాక్యాన్ని}{!!a{sentence}} (లేదా \tetoken{వాక్యాలు}{!!{sentence}s} సమితిని) సత్యం చేసే వాటిగా ఏ
\tetoken{నిర్మాణాలు}{!!{structure}s} సముదాయాలను వర్ణించగలం?
\end{explain}

\begin{defn}[ఒక సమితికి నమూనా]
$\Gamma$, భాష~$\Lang L$లోని \tetoken{వాక్యాలు}{!!{sentence}s} సమితి అనుకుందాం. ప్రతి
$!A \in \Gamma$కు $\Sat{M}{!A}$ అయితే, \tetoken{ఒక నిర్మాణం}{!!a{structure}}~$\Struct M$ను
$\Gamma$కు \emph{ఒక నమూనా} అంటాం.
\end{defn}

\begin{ex}
వాక్యం~$\lforall[x][x \le x]$, $\Assign{\le}{M}$ స్వావర్తన సంబంధం
అయినప్పుడు మాత్రమే~$\Struct M$లో సత్యం. వాక్యం
$\lforall[x][\lforall[y][((x \le y \land y \le x) \lif x = y)]]$,
$\Assign{\le}{M}$ ప్రతిసౌష్ఠవ సంబంధం అయినప్పుడు మాత్రమే~$\Struct M$లో
సత్యం. వాక్యం $\lforall[x][\lforall[y][\lforall[z][((x \le y \land y \le z)
      \lif x \le z)]]]$, $\Assign{\le}{M}$ సంక్రామక సంబంధం అయినప్పుడు
మాత్రమే~$\Struct M$లో సత్యం. అందువల్ల
\begin{align*}
\{\quad &\lforall[x][x \le x], \\
   & \lforall[x][\lforall[y][((x \le y \land y \le
    x) \lif x = y)]], \\
   &\lforall[x][\lforall[y][\lforall[z][((x \le y
      \land y \le z) \lif x \le z)]]] \quad \}
\end{align*}
యొక్క నమూనాలు సరిగ్గా~$\Assign{\le}{M}$ స్వావర్తన, ప్రతిసౌష్ఠవ,
సంక్రామక సంబంధం, అంటే పాక్షిక క్రమం, అయ్యే నిర్మాణాలే. కనుక వీటిని
\emph{పాక్షిక క్రమాల మొదటిస్థాయి సిద్ధాంతానికి} స్వీకృతాలుగా తీసుకోవచ్చు.
\end{ex}

\end{document}
